%-------------------------
% Resume in Latex
% Author : Pavel Kolyagin
% Based off of: https://github.com/sb2nov/resume
% License : MIT
%------------------------

\documentclass[letterpaper,12pt]{article}

\usepackage{latexsym}
\usepackage[empty]{fullpage}
\usepackage{titlesec}
\usepackage{marvosym}
\usepackage[usenames,dvipsnames]{color}
\usepackage{verbatim}
\usepackage{enumitem}
\usepackage[hidelinks]{hyperref}
\usepackage{fancyhdr}
\usepackage[T2A]{fontenc}
\usepackage[utf8]{inputenc}
\usepackage[english,russian]{babel}
\usepackage{tabularx}
\input{glyphtounicode}


%----------FONT OPTIONS----------
% sans-serif
% \usepackage[sfdefault]{FiraSans}
% \usepackage[sfdefault]{roboto}
% \usepackage[sfdefault]{noto-sans}
% \usepackage[default]{sourcesanspro}

% serif
% NB: PT Serif выбран вместо дефолтного Computer Modern — у CMR кернинг
% заглавных кириллических букв узкий, из-за чего "Колягин Павел" в \Huge
% визуально слипается и хуже извлекается в текст из PDF. PT Serif — T2A шрифт,
% разработанный под кириллицу, даёт заметно больший межсловный интервал.
\usepackage{PTSerif}
% \usepackage{CormorantGaramond}
% \usepackage{charter}


\pagestyle{fancy}
\fancyhf{} % clear all header and footer fields
\fancyfoot{}
\renewcommand{\headrulewidth}{0pt}
\renewcommand{\footrulewidth}{0pt}

% Adjust margins
\addtolength{\oddsidemargin}{-0.5in}
\addtolength{\evensidemargin}{-0.5in}
\addtolength{\textwidth}{1in}
\addtolength{\topmargin}{-.5in}
\addtolength{\textheight}{1.0in}

\urlstyle{same}

\raggedbottom
\raggedright
\setlength{\tabcolsep}{0in}

% Sections formatting
% NB: \scshape убран намеренно — кириллические псевдо-капсы плохо парсятся частью ATS.
\titleformat{\section}{
  \vspace{-4pt}\bfseries\raggedright\large
}{}{0em}{}[\color{black}\titlerule \vspace{-5pt}]

% Ensure that generate pdf is machine readable/ATS parsable
\pdfgentounicode=1

%-------------------------
% Custom commands
\newcommand{\resumeItem}[1]{
  \item\small{
    {#1 \vspace{-2pt}}
  }
}

\newcommand{\resumeSubheading}[4]{
  \vspace{-2pt}\item
    \begin{tabular*}{0.97\textwidth}[t]{l@{\extracolsep{\fill}}r}
      \textbf{#1} & #2 \\
      \textit{\small#3} & \textit{\small #4} \\
    \end{tabular*}\vspace{-7pt}
}

\newcommand{\resumeSubSubheading}[2]{
    \item
    \begin{tabular*}{0.97\textwidth}{l@{\extracolsep{\fill}}r}
      \textit{\small#1} & \textit{\small #2} \\
    \end{tabular*}\vspace{-7pt}
}

\newcommand{\resumeProjectHeading}[2]{
    \item
    \begin{tabular*}{0.97\textwidth}{l@{\extracolsep{\fill}}r}
      \small#1 & #2 \\
    \end{tabular*}\vspace{-7pt}
}

\newcommand{\resumeSubItem}[1]{\resumeItem{#1}\vspace{-4pt}}

\renewcommand\labelitemii{$\vcenter{\hbox{\tiny$\bullet$}}$}

\newcommand{\resumeSubHeadingListStart}{\begin{itemize}[leftmargin=0.15in, label={}]}
\newcommand{\resumeSubHeadingListEnd}{\end{itemize}}
\newcommand{\resumeItemListStart}{\begin{itemize}}
\newcommand{\resumeItemListEnd}{\end{itemize}\vspace{-5pt}}

%-------------------------------------------
%%%%%%  RESUME STARTS HERE  %%%%%%%%%%%%%%%%%%%%%%%%%%%%


\begin{document}

%----------HEADING----------
\begin{center}
    \textbf{\Huge Колягин Павел} \\ \vspace{2pt}
    \small Middle QA Automation / SDET $|$ Москва $|$ Удалённо \\ \vspace{1pt}
    \small +7 914-978-82-75 $|$ \href{mailto:inbox@paulrainy.ru}{\underline{inbox@paulrainy.ru}} $|$
    \href{https://github.com/paulrainy}{\underline{github.com/paulrainy}}
\end{center}


%-----------ОБО МНЕ-----------
\section{Обо мне}
  \small{Инженер по автоматизации тестирования с fullstack-профилем. Проектирую и автоматизирую API и UI тесты банковских систем и строю инструменты, ускоряющие сам процесс тестирования: аналитика нагрузочного тестирования, мониторинг стендов, RAG-генерация автотестов.}


%-----------ОПЫТ РАБОТЫ-----------
\section{Опыт работы}
  \resumeSubHeadingListStart

    \resumeSubheading
      {Инженер по автоматизации тестирования}{Июнь 2024 --- настоящее время}
      {ООО «РСХБ-Автоматизация» и ООО «РСХБ-ИНТЕХ»}{Россия, Москва}
      \resumeItemListStart
        \resumeItem{Автоматизирую регрессионное тестирование Единого фронтального решения (ЕФР) --- фронт-офисной платформы банка, экосистемы из 50+ микросервисов с параллельной миграцией фронтенда}
        \resumeItem{Спроектировал 86 API тест-кейсов и разработал 164 автоматизированных API- и UI-теста за второе полугодие 2025 (Java 21, TestNG, Selenide, Rest-assured, Allure; TMS --- Test IT; запуск в Jenkins)}
        \resumeItem{Внедрил слой SQL-валидации бизнес-данных в тест-фреймворк (Spring NamedParameterJdbcTemplate): 80+ параметризованных запросов с JOIN по кредитным, депозитным и ипотечным сущностям для сверки данных после операций}
        \resumeItem{Тестировал асинхронные интеграции через IBM MQ там, где REST API был недоступен --- извлечение уведомления из журнала ЕФР для верификации сценария создания клиента}
        \resumeItem{Построил RAG-пайплайн генерации API-автотестов из тест-кейсов (Python, Sentence-Transformers, Milvus, Apache Airflow): поиск дубликатов и связанных тест-кейсов ускорился с нескольких часов до 0.12 с (векторный поиск по Milvus), генерация одного API-автотеста --- 20 с против 30 минут вручную; 3 квартала пайплайн закрывал план 150 API-автотестов в квартал}
        \resumeItem{Устранил flaky-тесты, вызванные конкуренцией за общие данные: перевёл автотесты с единого захардкоженного тестового пользователя на генерацию отдельного пользователя под каждый запуск}
        \resumeItem{Выявил и отозвал дефектный тест-кейс (смена прав администратора в ЕФР без отката при падении), каскадно ронявший сотни автотестов в одном прогоне Jenkins --- усилил требования к изоляции состояния тестов}
        \resumeItem{Разработал инструмент анализа результатов нагрузочного тестирования (Python, Streamlit, PostgreSQL): интеграция с 6 источниками метрик (Zabbix, Prometheus, InfluxDB, Grafana, Kibana, PostgreSQL), сравнение до 5 релизов одновременно, сокращение сбора аналитики с 3--8 часов до 2 минут, экспорт отчётов в web, docx и xlsx}
        \resumeItem{Покрыл инструмент анализа результатов нагрузочного тестировани 593 автотестами (unit / integration / e2e на pytest + Playwright, покрытие 72\%, GitLab CI)}
        \resumeItem{Разработал EFR Down Detector --- сервис мониторинга (Python, FastAPI, Redis, Vue 3): в реальном времени отслеживает доступность REST-сервисов ЕФР на 4 тестовых окружениях, которым пользуются 14 тестировщиков вместо ручной проверки стендов}
        \resumeItem{Провёл рефакторинг генератора тестовых клиентов банка: вынес логику в отдельный Spring Boot REST-сервис (Java 21) и Vue 3-фронтенд, контейнеризовал (Docker Compose), упростив поддержку и масштабирование}
        \resumeItem{Написал парсер для legacy почтового сервиса Lotus Notes (Java 21, Selenide): автоматизировал миграцию писем пользователей на новый почтовый сервис с сохранением вложений}
        \resumeItem{Интегрировал в корпоративный генератор презентаций и промо-материалов для HR-отдела поддержку новых LLM-моделей генерации изображений (Nano Banana, GPT-5-image)}
        \resumeItem{Исследовал embedding- и reranker-модели для актуализации пар «тест-кейс~$\leftrightarrow$~автотест» (Jina embeddings, Test IT API)}
      \resumeItemListEnd

  \resumeSubHeadingListEnd


%-----------НАГРАДЫ И ВЫСТУПЛЕНИЯ-----------
\section{Награды и выступления}
  \resumeSubHeadingListStart
    \resumeProjectHeading
      {\textbf{Победитель} $|$ \emph{«Лучший по профессии», номинация «Лучший тестировщик» (РСХБ-Авто)}}{05.2026}
    \resumeProjectHeading
      {\textbf{Спикер} $|$ \emph{Testify \#12, QA-митап}}{12.2025}
    \resumeProjectHeading
      {\textbf{Стендист} $|$ \emph{Heisenbug Spring 2025, конференция по тестированию ПО}}{04.2025}
  \resumeSubHeadingListEnd


%-----------ОБРАЗОВАНИЕ-----------
\section{Образование}
  \resumeSubHeadingListStart
    \resumeSubheading
      {РТУ МИРЭА (Российский технологический университет «МИРЭА»)}{Россия, Москва}
      {Магистр, Информатика и вычислительная техника}{Сентябрь 2023 --- Июль 2025}
    \resumeSubheading
      {НИУ МЭИ (Национальный исследовательский университет «МЭИ»)}{Россия, Москва}
      {Бакалавр, Информационная безопасность}{Сентябрь 2019 --- Июль 2023}
  \resumeSubHeadingListEnd


%-----------ДОПОЛНИТЕЛЬНОЕ ОБРАЗОВАНИЕ-----------
\section{Дополнительное образование}
  \resumeSubHeadingListStart
    \resumeSubheading
      {Академия трансформации (АО «Россельхозбанк»)}{Россия, Москва}
      {Слушатель}{Апрель 2026}
      \resumeItemListStart
        \resumeItem{«Оптимизация процессов. Подход DMAIC» --- факультет «Управление процессами»}
      \resumeItemListEnd
    \resumeSubheading
      {РТУ МИРЭА (Российский технологический университет «МИРЭА»)}{Россия, Москва}
      {Студент}{Сентябрь 2023 --- Декабрь 2023}
      \resumeItemListStart
        \resumeItem{Менеджмент (профессиональная переподготовка)}
      \resumeItemListEnd
  \resumeSubHeadingListEnd


%-----------СТЕК ТЕХНОЛОГИЙ-----------
\section{Стек технологий}
 \begin{itemize}[leftmargin=0.15in, label={}]
    \small{\item{
     \textbf{Testing}{: REST API testing, UI-автоматизация, pytest, Playwright, Selenide, Rest-assured, Allure, Test IT} \\
     \textbf{Programming}{: Java, Python} \\
     \textbf{Backend}{: Spring, FastAPI, SQLAlchemy} \\
     \textbf{Frontend}{: Vue 3, Streamlit} \\
     \textbf{Databases}{: PostgreSQL, Milvus, QDrant, Redis} \\
     \textbf{CI/CD \& DevOps}{: Git, Docker, Docker-compose, Jenkins, GitLab CI, GitHub Actions, Apache Airflow} \\
     \textbf{AI \& ML}{: RAG, embeddings \& rerankers, Sentence-Transformers, camel-ai, pydantic-ai} \\
     \textbf{Языки}{: Русский (родной), Английский (B1)}
    }}
 \end{itemize}


%-------------------------------------------
\end{document}
